% !TEX root = ../main.tex
%
\pdfbookmark[0]{Abstract}{Abstract}
\chapter*{Abstract}
\label{sec:abstract}
\vspace*{-10mm}

Standard uncertainty estimation in Large Language Models (LLMs) 
relies on single, fixed probability distributions. This 
``point-estimate'' approach often conflates benign linguistic 
diversity with factual ignorance, failing to reliably quantify 
uncertainty.
In this thesis, we introduce \textit{Credal Semantic Entropy},
a rigorous framework grounded in Imprecise Probabilities. 
We propose a method to construct Credal Sets---convex sets of 
probability distributions---based on the statistical notion of 
relative likelihood. Instead of relying on a single temperature,
we identify a set of statistically valid temperatures that
maximize the likelihood of calibration data. We cluster generated
sequences by semantic meaning and derive interval-valued
probabilities, computing the Upper Semantic Entropy to capture
a worst-case estimate of uncertainty.
We empirically evaluate our approach on the CoQA benchmark using
the OPT-6.7B model. While theoretically robust, experiments 
reveal that the credal framework is highly sensitive to relative
likelihood thresholds ($\alpha$-cuts). Under our tested 
configurations, single-point baselines outperform Credal Semantic 
Entropy in uncertainty estimation (measured via AUROC).
%This thesis lays the foundational groundwork for
%applying credal uncertainty estimation to autoregressive language
%models, opening new avenues for future research.

\iffalse
Accurately quantifying uncertainty in Large Language Models (LLMs) is 
crucial for their deployment in safety-critical domains. However, 
standard uncertainty estimation methods typically rely on a single 
probability distribution defined by fixed decoding hyperparameters. 
We argue that this "point-estimate" approach fails to adequately 
capture epistemic uncertainty, often conflating benign linguistic 
diversity with factual ignorance.
In this thesis, we introduce \textit{Credal Semantic Entropy}, a 
rigorous framework for hallucination detection grounded in the theory 
of Imprecise Probabilities. We propose a method to construct Credal 
Sets—convex sets of probability distributions—based on the statistical 
notion of relative likelihood. Instead of relying on a single 
temperature, we identify a set of statistically valid temperatures 
that maximize the likelihood of calibration data. We then aggregate 
generated sequences into semantic clusters and derive interval-valued 
probabilities  for each meaning. By computing the Upper Semantic 
Entropy over these intervals, we obtain a worst-case estimate of 
uncertainty.
\fi

\vspace*{20mm}
